% !TEX program = xelatex
\documentclass[11pt,a4paper]{article}
\usepackage{fontspec}
\setmainfont{Calibri}
\setmonofont[Scale=0.88]{Consolas}
\newfontfamily\symbolfont{Segoe UI Symbol}
\usepackage{polyglossia}\setdefaultlanguage{polish}
\usepackage[table]{xcolor}
\usepackage{booktabs,longtable,array,geometry,enumitem,graphicx,fancyhdr,caption}
\usepackage[most]{tcolorbox}
\PassOptionsToPackage{hyphens}{url}
\usepackage[hidelinks,unicode]{hyperref}
\emergencystretch=2.5em
\geometry{a4paper,margin=2.1cm,top=2.3cm,bottom=2.3cm}
\setlength{\parindent}{0pt}\setlength{\parskip}{5pt}
\definecolor{apteal}{HTML}{0E7C86}
\definecolor{apteallight}{HTML}{E6F4F5}
\definecolor{apgrey}{HTML}{F3F4F6}
\definecolor{apamber}{HTML}{B45309}
\definecolor{apamberlight}{HTML}{FEF3C7}
\definecolor{apred}{HTML}{B91C1C}
\definecolor{apredlight}{HTML}{FEE2E2}
\definecolor{apgreen}{HTML}{15803D}
\definecolor{apgreenlight}{HTML}{DCFCE7}
\renewcommand{\arraystretch}{1.25}
\captionsetup{font=small,labelfont=bf}
\pagestyle{fancy}\fancyhf{}
\fancyfoot[L]{\small AMPERE POINT — plan testów offline · Q11 z modułem Shelly · v3 · 2026-09-28}
\fancyfoot[R]{\small\thepage}
\renewcommand{\headrulewidth}{0pt}
\newtcolorbox{uwaga}[1][Uwaga]{enhanced,breakable,colback=apamberlight,colframe=apamber,title={#1},fonttitle=\bfseries,boxrule=0.6pt,arc=2pt,left=6pt,right=6pt,top=4pt,bottom=4pt}
\newtcolorbox{info}[1][Jak to działa]{enhanced,breakable,colback=apteallight,colframe=apteal,title={#1},fonttitle=\bfseries,boxrule=0.6pt,arc=2pt,left=6pt,right=6pt,top=4pt,bottom=4pt}
\newtcolorbox{wazne}[1][Ważne]{enhanced,breakable,colback=apredlight,colframe=apred,title={#1},fonttitle=\bfseries,boxrule=0.6pt,arc=2pt,left=6pt,right=6pt,top=4pt,bottom=4pt}
\newtcolorbox{przyklad}[1][Przykład liczbowy]{enhanced,breakable,colback=apgreenlight,colframe=apgreen,title={#1},fonttitle=\bfseries,boxrule=0.6pt,arc=2pt,left=6pt,right=6pt,top=4pt,bottom=4pt}
\newtcolorbox{kodbox}{enhanced,breakable,colback=apgrey,colframe=apgrey,boxrule=0pt,arc=2pt,left=5pt,right=5pt,top=3pt,bottom=3pt,fontupper=\ttfamily\footnotesize}
\setlist{nosep,leftmargin=*}
\setcounter{secnumdepth}{2}
\begin{document}

{\LARGE\bfseries Plan testów offline: Q11 z modułem Shelly\par}
{\large Co jeszcze trzeba sprawdzić, żeby potwierdzić, że bez chmury i bez internetu Q11 na Shelly umie co najmniej to, co Q11 na Tuya, że da się nią sterować prosto z telefonu, i że po zaniku zasilania pamięta to, co powinna · AMPERE POINT · 28 września 2026 · oprac. Dawid Cekała\par}
\vspace{4pt}\textcolor{apteal}{\rule{\textwidth}{1.2pt}}

\section*{Co zmieniło się w v3}

\begin{longtable}{>{\raggedright\arraybackslash}p{1.65cm}>{\raggedright\arraybackslash}p{14.07cm}}
\toprule
\rowcolor{apteallight}\textbf{Miejsce} & \textbf{Zmiana} \\
\midrule\endhead
\bottomrule\endfoot
Rozdział 3 & Trzy tabele porównania Tuya i Shelly: bez chmury, bez internetu, bez Wi-Fi. Przy każdej funkcji osobno: czy działa i czy to sprawdziliśmy \\
T1 & Nowy krok 1g: aplikacja Shelly w sieci lokalnej, moduł bez internetu \\
T2 & Nowy krok 2f: MQTT i most OCPP bez internetu \\
T7 & Nowy krok 7g: Home Assistant z oficjalną integracją Shelly \\
\end{longtable}

Wersja v2 wprowadziła poziomy łączności, słowniczek, test API i poprawki po przeglądzie v1; zostaje obok jako plik archiwalny.


\section*{W skrócie}

\begin{itemize}
\item \textbf{Cel dla użytkownika.} Zwykły użytkownik, bez Home Assistant, ma prosto sterować ładowarką bez chmury: z aplikacji przez Bluetooth albo ze strony modułu w sieci lokalnej.
\item \textbf{Co Tuya daje bez chmury.} Większość funkcji, które działają bez internetu, siedzi w sterowniku Q11: menu, przyciski, start po podłączeniu auta, limit prądu, okno godzin, limit energii. Sterowanie z daleka idzie przez chmurę Tuya. Harmonogramy z aplikacji zapisuje chmura; kto je wykonuje, chmura czy moduł Tuya, jeszcze nie wiadomo. Lokalnie Tuyą da się sterować kluczem lokalnym przez Home Assistant, ale to droga dla zaawansowanych. Sama aplikacja Tuya, gdy telefon jest w tej samej sieci co ładowarka, też wysyła polecenia lokalnie. Nie sprawdzaliśmy, czy robi to także bez internetu, skoro do logowania potrzebuje chmury.
\item \textbf{Co sprawdzamy.} Trzy rzeczy. Czy funkcje sterownika działają z modułem Shelly bez internetu. Czy to, co moduł Shelly daje od siebie bez chmury, naprawdę działa: Bluetooth, strona WWW, harmonogramy w module, webhooki, API. I co się dzieje po zaniku zasilania, kiedy dwa układy z własną pamięcią startują naraz.
\item \textbf{Kolejność od bezpiecznych do ryzykownych.} Reset sieci z menu ładowarki zostawiamy na koniec, bo może skasować Wi-Fi modułu.
\end{itemize}

\section{Poziomy łączności}

„Offline” może znaczyć kilka różnych rzeczy. W tym planie rozróżniamy pięć poziomów:


\begin{longtable}{>{\raggedright\arraybackslash}p{3.00cm}>{\raggedright\arraybackslash}p{4.28cm}>{\raggedright\arraybackslash}p{2.25cm}>{\raggedright\arraybackslash}p{5.31cm}}
\toprule
\rowcolor{apteallight}\textbf{Poziom} & \textbf{Co działa} & \textbf{Czego nie ma} & \textbf{Przykład z życia} \\
\midrule\endhead
\bottomrule\endfoot
A. Z chmurą & Wi-Fi, internet, chmura producenta & — & zwykła praca; aplikacja z dowolnego miejsca \\
B. Bez chmury & Wi-Fi i internet & chmura producenta wyłączona w module & nasze próby od 25.09 \\
C. Bez internetu & Wi-Fi i sieć lokalna & internet & awaria łącza; firma, która nie wpuszcza ładowarek do internetu \\
D. Bez Wi-Fi & Bluetooth albo własna sieć modułu, telefon obok ładowarki & router & garaż bez zasięgu Wi-Fi \\
E. Bez niczego & sama ładowarka & żadnej łączności & tylko menu i przyciski \\
\end{longtable}

„Offline” w tym planie oznacza poziomy C, D i E. Na poziomie B większość już sprawdziliśmy (tabela 3.1). Brak internetu symulowaliśmy dotąd tylko dla serwera czasu.


\begin{longtable}{>{\raggedright\arraybackslash}p{12.35cm}>{\raggedright\arraybackslash}p{3.37cm}}
\toprule
\rowcolor{apteallight}\textbf{Test} & \textbf{Poziom} \\
\midrule\endhead
\bottomrule\endfoot
T1 Bluetooth, strona WWW przez sieć modułu, aplikacja & D; krok 1g na C \\
T2 funkcje sterownika & C i E \\
T3 okno godzin, tryb, limit energii & C i E \\
T4 harmonogram, webhook, pamięć w module & C \\
T5 zanik zasilania & C \\
T6 reset sieci z menu & E, skutek dla modułu \\
T7 API & C \\
\end{longtable}

\section*{Słowniczek}

\begin{itemize}
\item \textbf{Menu ładowarki.} Ustawienia na samej Q11: wyświetlacz i dwa przyciski na obudowie, 5 języków. Działa bez żadnego modułu. W planie „z menu” znaczy: ktoś stoi przy ładowarce i ustawia przyciskami.
\item \textbf{Okno godzin ładowania.} Ustawienie sterownika Q11: godzina od i godzina do, w których wolno ładować, np. 22–6 dla nocnej taryfy. Działa w trybie „harmonogram”. Sterownik pilnuje okna sam, potrzebuje tylko aktualnej godziny od modułu. Zapisane w dwóch bajtach: godzina startu i godzina końca, \textbf{bez minut}.
\item \textbf{Harmonogram w aplikacji Tuya.} Wpisy tworzy aplikacja, a przechowuje chmura. Telefon do wykonania nie jest potrzebny. Kto wykonuje wpis, nie jest rozstrzygnięte, bo źródła są sprzeczne. Za chmurą przemawia definicja produktu Q21 w Tuya, w której stoi \emph{CloudTiming: „cloud timing without local timing”}. Za modułem Tuya przemawiają logi ładowarki testowej (Q11 PRO): wykonanie wpisu widać tylko jako raport samej ładowarki, bez polecenia z chmury, które pojawia się przy scenach. Rozstrzyga test T0b.
\item \textbf{Harmonogram w module Shelly.} Do 20 zadań zapisanych w samym module, np. „o 22:00 ustaw 16 A”. Moduł wykonuje je sam, bez chmury i internetu. Potrzebuje godziny, więc przy braku internetu wymaga lokalnego serwera czasu.
\item \textbf{API.} Zestaw poleceń, którymi inne programy rozmawiają z modułem, np. „podaj stan” albo „ustaw limit 10 A”. Te same polecenia idą różnymi drogami: HTTP, WebSocket, MQTT, Bluetooth.
\item \textbf{Webhook.} Adres WWW, który moduł \textbf{sam} wywołuje, gdy coś się stanie, np. „zaczęło się ładowanie”. To odwrotność API: przy API program pyta moduł, przy webhooku moduł sam zawiadamia program. Przykład: moduł wywołuje \texttt{http:\allowbreak{}/\allowbreak{}/\allowbreak{}192.\allowbreak{}168.\allowbreak{}0.\allowbreak{}57:\allowbreak{}8080/\allowbreak{}zdarzenie}, a program na laptopie zapisuje zdarzenie.
\item \textbf{Zapamiętywanie pola.} Ustawienie pola w portalu Shelly („persisted”): moduł trzyma wartość w swojej pamięci i po restarcie ją odtwarza. U nas zapamiętywane są limit prądu i włącznik ładowania.
\item \textbf{Sieć modułu.} Moduł może sam wystawić sieć Wi-Fi (punkt dostępowy), do której łączy się telefon, bez routera. \textbf{Stacja Wi-Fi} to druga strona: moduł jako klient sieci biura.
\end{itemize}

\section{Tuya a Shelly na każdym poziomie łączności}

Trzy tabele, po jednej na poziom. W każdej: co może Q11 na Tuya, co może Q11 na Shelly i czy to sprawdziliśmy.


\textbf{Czy działa:} \leavevmode{\symbolfont\color{apgreen}\textbf{✔}} działa · \leavevmode{\symbolfont\color{apamber}\textbf{◐}} częściowo · \leavevmode{\symbolfont\color{apred}\textbf{✘}} nie działa · ? nie wiadomo.


\textbf{Sprawdzone?:}

\begin{itemize}
\item \textbf{tak} — sprawdzone u nas na sprzęcie, z datą;
\item \textbf{z użytkowania} — znane z pracy ładowarek u klientów i w serwisie, bez osobnego testu;
\item \textbf{z dokumentacji} — wiemy od producenta, nie sprawdzaliśmy;
\item \textbf{wniosek} — wynika z budowy systemu, nikt nie sprawdzał;
\item \textbf{nie {\symbolfont→} T…} — sprawdzi test z tego planu.
\end{itemize}

Funkcje samego sterownika, czyli menu, przyciski, start po podłączeniu auta i limit z menu, nie zależą od łączności. Są w tabeli 3.3, ale obowiązują na każdym poziomie.


\subsection{Bez chmury: Wi-Fi i internet są, chmury producenta nie ma}

Tuya nie ma wyłącznika chmury. Dla Tuya ten poziom znaczy: internet jest, ale z chmury Tuya nie korzystamy, bo jest niedostępna albo zablokowana.


\begin{longtable}{>{\raggedright\arraybackslash}p{3.00cm}>{\raggedright\arraybackslash}p{3.47cm}>{\raggedright\arraybackslash}p{2.34cm}>{\raggedright\arraybackslash}p{3.25cm}>{\raggedright\arraybackslash}p{2.34cm}}
\toprule
\rowcolor{apteallight}\textbf{Funkcja} & \textbf{Q11 na Tuya} & \textbf{Sprawdzone?} & \textbf{Q11 na Shelly} & \textbf{Sprawdzone?} \\
\midrule\endhead
\bottomrule\endfoot
Sterowanie z telefonu spoza domu & \leavevmode{\symbolfont\color{apred}\textbf{✘}} zdalny dostęp idzie przez chmurę & wniosek & \leavevmode{\symbolfont\color{apred}\textbf{✘}} tak samo & wniosek \\
Aplikacja producenta w tej samej sieci & ? polecenia idą lokalnie, ale do logowania aplikacja potrzebuje chmury & polecenia lokalne: tak, 07.09; bez chmury: nie & ? & nie {\symbolfont→} T1g \\
Strona WWW modułu w przeglądarce & \leavevmode{\symbolfont\color{apred}\textbf{✘}} moduł Tuya jej nie ma & z dokumentacji & \leavevmode{\symbolfont\color{apamber}\textbf{◐}} działa; na produkcie Q11 jeszcze nie sprawdzona & tak, 24.09 na symulatorze; Q11 {\symbolfont→} T1f \\
Sterowanie z komputera w sieci lokalnej & \leavevmode{\symbolfont\color{apamber}\textbf{◐}} tylko kluczem lokalnym, który trzeba raz pobrać z konta Tuya & tak, nasza integracja Home Assistant & \leavevmode{\symbolfont\color{apgreen}\textbf{✔}} HTTP i WebSocket, bez klucza & tak, 25.09 \\
Home Assistant & \leavevmode{\symbolfont\color{apamber}\textbf{◐}} tuya local z naszym profilem & tak, 24.09 & ? oficjalna integracja Shelly & nie {\symbolfont→} T7g \\
MQTT do własnego systemu & \leavevmode{\symbolfont\color{apred}\textbf{✘}} moduł łączy się tylko z brokerem w chmurze Tuya & wniosek & \leavevmode{\symbolfont\color{apgreen}\textbf{✔}} własny broker & tak, 25.09 \\
OCPP & \leavevmode{\symbolfont\color{apred}\textbf{✘}} & wniosek & \leavevmode{\symbolfont\color{apamber}\textbf{◐}} przez most; bez pełnej transakcji & tak, 25.09 \\
Godzina dla sterownika & ? przychodzi z chmury Tuya & wniosek & \leavevmode{\symbolfont\color{apgreen}\textbf{✔}} z publicznego serwera czasu & tak, 25.09, 11:51 \\
Harmonogram & ? zależy, kto wykonuje wpis: chmura czy moduł & nie {\symbolfont→} T0b & ? do 20 zadań w module & nie {\symbolfont→} T4 \\
Powiadomienie do własnego systemu & \leavevmode{\symbolfont\color{apred}\textbf{✘}} & wniosek & ? webhook & nie {\symbolfont→} T4 \\
Ikona sieci na wyświetlaczu & ? & nie & \leavevmode{\symbolfont\color{apamber}\textbf{◐}} pokazuje „chmura”, choć chmury nie ma & tak, 25.09 \\
Aktualizacja oprogramowania modułu & \leavevmode{\symbolfont\color{apred}\textbf{✘}} tylko z chmury Tuya & wniosek & \leavevmode{\symbolfont\color{apamber}\textbf{◐}} z portalu przez Bluetooth; komputer potrzebuje internetu, moduł nie & tak, 24.09 \\
\end{longtable}

\textbf{Wniosek.} Shelly bez chmury traci zdalny dostęp z telefonu, a zyskuje wszystko lokalnie: API, MQTT, OCPP przez most i godzinę z internetu. Większość z tego sprawdziliśmy 24–25.09. Tuya bez chmury działa lokalnie tylko u kogoś, kto pobierze klucz lokalny i postawi Home Assistant.


\subsection{Bez internetu: Wi-Fi i sieć lokalna są, internetu nie ma}

Zdalny dostęp i wszystko, co idzie przez chmurę, odpada w obu systemach, więc tych wierszy tu nie powtarzamy.


\begin{longtable}{>{\raggedright\arraybackslash}p{3.54cm}>{\raggedright\arraybackslash}p{2.52cm}>{\raggedright\arraybackslash}p{2.35cm}>{\raggedright\arraybackslash}p{2.94cm}>{\raggedright\arraybackslash}p{3.05cm}}
\toprule
\rowcolor{apteallight}\textbf{Funkcja} & \textbf{Q11 na Tuya} & \textbf{Sprawdzone?} & \textbf{Q11 na Shelly} & \textbf{Sprawdzone?} \\
\midrule\endhead
\bottomrule\endfoot
Aplikacja producenta w tej samej sieci & ? polecenia idą lokalnie, ale logowanie wymaga chmury & nie {\symbolfont→} T0a & ? & nie {\symbolfont→} T1g \\
Strona WWW modułu & \leavevmode{\symbolfont\color{apred}\textbf{✘}} & z dokumentacji & \leavevmode{\symbolfont\color{apamber}\textbf{◐}} działa; Q11 jeszcze nie & tak, 24.09 na symulatorze; Q11 {\symbolfont→} T1f \\
Sterowanie z komputera w sieci lokalnej & \leavevmode{\symbolfont\color{apamber}\textbf{◐}} kluczem lokalnym; ruch zostaje w sieci & nie & \leavevmode{\symbolfont\color{apgreen}\textbf{✔}} ruch zostaje w sieci & tak, 25.09, ale z internetem w tle; bez {\symbolfont→} T7 \\
Home Assistant & \leavevmode{\symbolfont\color{apamber}\textbf{◐}} tuya local & nie & ? & nie {\symbolfont→} T7g \\
MQTT, własny broker & \leavevmode{\symbolfont\color{apred}\textbf{✘}} & wniosek & \leavevmode{\symbolfont\color{apgreen}\textbf{✔}} & tak, 25.09, z internetem w tle; bez {\symbolfont→} T2f \\
OCPP przez most & \leavevmode{\symbolfont\color{apred}\textbf{✘}} & wniosek & \leavevmode{\symbolfont\color{apamber}\textbf{◐}} most w sieci lokalnej & tak, 25.09, z internetem w tle; bez {\symbolfont→} T2f \\
Godzina dla sterownika & \leavevmode{\symbolfont\color{apred}\textbf{✘}} czas tylko z chmury Tuya & wniosek & \leavevmode{\symbolfont\color{apamber}\textbf{◐}} tylko z lokalnym serwerem czasu; bez niego zerowa data & tak, 25.09, oba przypadki \\
Okno godzin ładowania & ? sterownik ma okno, ale nie ma godziny & nie & ? z lokalnym serwerem czasu & nie {\symbolfont→} T3 \\
Harmonogram & ? & nie {\symbolfont→} T0b & ? z lokalnym serwerem czasu & nie {\symbolfont→} T4 \\
Powiadomienie do własnego systemu & \leavevmode{\symbolfont\color{apred}\textbf{✘}} & wniosek & ? webhook & nie {\symbolfont→} T4 \\
Ikona sieci na wyświetlaczu & ? & nie & \leavevmode{\symbolfont\color{apamber}\textbf{◐}} pokazuje „chmura” & nie {\symbolfont→} T2e \\
Pamięć ustawień po zaniku zasilania & \leavevmode{\symbolfont\color{apamber}\textbf{◐}} sterownik pamięta swoje & wniosek & ? dwa układy z pamięcią & nie {\symbolfont→} T5 \\
\end{longtable}

\textbf{Wniosek.} W sieci lokalnej Shelly robi to samo co bez chmury. Jedyna różnica to godzina: trzeba lokalnego serwera czasu, a w przyszłości, jeśli Shelly to zmieni, wystarczy czas z telefonu. Tuya bez internetu to Home Assistant z kluczem lokalnym, w dodatku bez godziny dla sterownika.


\subsection{Bez Wi-Fi: tylko telefon albo laptop obok ładowarki}

\begin{longtable}{>{\raggedright\arraybackslash}p{3.73cm}>{\raggedright\arraybackslash}p{2.56cm}>{\raggedright\arraybackslash}p{2.35cm}>{\raggedright\arraybackslash}p{3.23cm}>{\raggedright\arraybackslash}p{2.52cm}}
\toprule
\rowcolor{apteallight}\textbf{Funkcja} & \textbf{Q11 na Tuya} & \textbf{Sprawdzone?} & \textbf{Q11 na Shelly} & \textbf{Sprawdzone?} \\
\midrule\endhead
\bottomrule\endfoot
Aplikacja producenta przez Bluetooth & ? moduł Tuya ma Bluetooth, sterowanie nim niesprawdzone & nie & ? aplikacja wymaga konta w chmurze & nie {\symbolfont→} T1c \\
Strona WWW przez sieć modułu & \leavevmode{\symbolfont\color{apred}\textbf{✘}} & z dokumentacji & \leavevmode{\symbolfont\color{apamber}\textbf{◐}} działa bez routera i bez aplikacji & tak, 24.09 na symulatorze; Q11 {\symbolfont→} T1f, T1d \\
Polecenia z komputera przez Bluetooth & \leavevmode{\symbolfont\color{apred}\textbf{✘}} brak otwartego sposobu & wniosek & ? włączone w konfiguracji modułu & konfiguracja: tak, 25.09; działanie {\symbolfont→} T1b \\
Godzina dla sterownika & \leavevmode{\symbolfont\color{apred}\textbf{✘}} & wniosek & \leavevmode{\symbolfont\color{apred}\textbf{✘}} czas ustawiony ręcznie nie trafia do sterownika & tak, 25.09 \\
Okno godzin ładowania & ? bez godziny & nie & ? bez godziny & nie {\symbolfont→} T3b \\
Harmonogram & \leavevmode{\symbolfont\color{apred}\textbf{✘}} bez sieci & wniosek & \leavevmode{\symbolfont\color{apamber}\textbf{◐}} działa, dopóki moduł pamięta godzinę sprzed utraty sieci; po zaniku zasilania nie & wniosek \\
MQTT, OCPP, Home Assistant & \leavevmode{\symbolfont\color{apred}\textbf{✘}} brak sieci & wniosek & \leavevmode{\symbolfont\color{apred}\textbf{✘}} brak sieci & wniosek \\
Start ładowania po podłączeniu auta & \leavevmode{\symbolfont\color{apgreen}\textbf{✔}} & z użytkowania & \leavevmode{\symbolfont\color{apgreen}\textbf{✔}} & tak, 25.09 \\
Limit prądu z menu ładowarki & \leavevmode{\symbolfont\color{apgreen}\textbf{✔}} & z użytkowania & ? & nie {\symbolfont→} T2b \\
Tryb, okno godzin i limit energii z menu & \leavevmode{\symbolfont\color{apgreen}\textbf{✔}} ustawienie działa; okno wymaga godziny & z użytkowania & ? & nie {\symbolfont→} T3 \\
Reset sieci z menu ładowarki & \leavevmode{\symbolfont\color{apgreen}\textbf{✔}} moduł wchodzi w parowanie & z dokumentacji Tuya & ? & nie {\symbolfont→} T6 \\
\end{longtable}

\textbf{Wniosek.} Bez Wi-Fi Tuya ma tylko menu ładowarki. Shelly ma dodatkowo stronę WWW przez własną sieć modułu, bez aplikacji i bez konta; na symulatorze to działało. Przeszkodą jest godzina: bez Wi-Fi nie ma serwera czasu, a czasu ustawionego ręcznie moduł sterownikowi nie przekazuje. Dopóki Shelly tego nie zmieni, okno godzin i harmonogramy bez Wi-Fi nie mają godziny po zaniku zasilania.


\textbf{Porównanie z Tuya, opcjonalnie (T0).} Dwie niewiadome po stronie Tuya da się rozstrzygnąć tylko na Q11 z modułem Tuya odciętym od internetu:

\begin{itemize}
\item \textbf{T0a:} czy aplikacja Tuya steruje przez sieć lokalną.
\item \textbf{T0b:} czy wpis harmonogramu wykona się bez internetu; jeśli tak, wykonuje go moduł Tuya, a nie chmura.
\end{itemize}

Odcięcie jednej ładowarki od internetu wymaga reguły w routerze, którą ustawia administrator sieci. Robimy to tylko wtedy, jeśli będzie taka możliwość. Sam telefon bez sieci niczego tu nie rozstrzyga, bo ładowarka nadal ma internet.


\section{Założenia i przygotowanie}

\begin{itemize}
\item Chmura Shelly wyłączona przez cały czas, jak dotąd.
\item W sieci działają kontenery: broker MQTT i lokalny serwer czasu (192.168.0.57). Serwer czasu musi działać \textbf{przed} każdym włączeniem ładowarki, bo moduł pyta o czas zaraz po starcie.
\item Monitor logu przez Wi-Fi uruchomiony przed każdym testem. Przy zaniku zasilania monitor sam się wznawia po powrocie modułu.
\item Symulator auta podłączony: stany „podłączone” (9 V) i „ładuje” (6 V). Bez obciążenia prądowego, więc prąd i moc faz zostają poza tym planem.
\item Laptop ma kartę Bluetooth (Intel). Do testów z komputera służy biblioteka \texttt{bleak} w Pythonie. Telefon z aplikacją Shelly jest potrzebny do testu z użytkownikiem.
\item Wi-Fi laptopa i router zostają bez zmian.
\item \textbf{Moduł na poziomie C.} Moduł dostaje stały adres (ten sam, 192.168.0.238) bez bramy do internetu. Widzi laptopa, broker i serwer czasu, ale nie internet. Zmiana jest odwracalna z sieci lokalnej: powrót na adres automatyczny. Przed zmianą sprawdzamy, że sieć modułu jest włączona, bo to droga ratunkowa, gdyby moduł zniknął z sieci.
\item Po każdym teście zapis w tabeli wyników: \leavevmode{\symbolfont\color{apgreen}\textbf{✔}} działa, \leavevmode{\symbolfont\color{apamber}\textbf{◐}} działa z zastrzeżeniem, \leavevmode{\symbolfont\color{apred}\textbf{✘}} nie działa, z opisem tego, co zaobserwowano.
\end{itemize}

\section{Testy}

\subsection{T1. Bluetooth i strona WWW}

\textbf{Cel.} Potwierdzić, że ładowarką da się sterować bez Wi-Fi i bez internetu, z telefonu i z komputera.


\begin{longtable}{>{\raggedright\arraybackslash}p{3.00cm}>{\raggedright\arraybackslash}p{6.14cm}>{\raggedright\arraybackslash}p{6.14cm}}
\toprule
\rowcolor{apteallight}\textbf{Krok} & \textbf{Co robimy} & \textbf{Wynik oczekiwany} \\
\midrule\endhead
\bottomrule\endfoot
1a & Skan Bluetooth z laptopa & moduł widoczny pod swoją nazwą, z reklamą Shelly \\
1b & Polecenie przez Bluetooth z laptopa: odczyt stanu, potem zmiana limitu prądu na 10 A & sterownik potwierdza w logu w ciągu 1 s, tak jak przy Wi-Fi \\
1c & Telefon z wyłączonym Wi-Fi i danymi komórkowymi, aplikacja Shelly przez Bluetooth: włącz/wyłącz ładowanie, zmiana limitu & \textbf{do ustalenia.} Aplikacja wymaga konta w chmurze, a chmura modułu jest wyłączona; może nie pokazać urządzenia. Zapisujemy, co widać i ile trwa połączenie \\
1f & Telefon łączy się z siecią modułu, przeglądarka otwiera stronę modułu, zmiana limitu i włącznika & działa bez routera i bez aplikacji; 24.09 działało na symulatorze, teraz sprawdzamy na produkcie Q11 \\
1d & Moduł bez Wi-Fi: przez Bluetooth wyłączamy stację Wi-Fi modułu. Wcześniej sprawdzamy, że sieć modułu jest włączona jako droga awaryjna, i \textbf{czy ma hasło}. Powtarzamy 1b, 1c i 1f & działa bez Wi-Fi; log niedostępny, więc dowodem są wyświetlacz ładowarki i odpowiedź przez Bluetooth \\
1g & Poziom C: telefon w sieci biura, aplikacja Shelly; moduł bez internetu i bez chmury & czy aplikacja znajduje moduł w sieci lokalnej i steruje nim. Telefon ma w tym czasie internet przez sieć biura, więc sprawdzamy tylko drogę telefon {\symbolfont→} moduł \\
1e & Przez Bluetooth włączamy stację Wi-Fi z powrotem & moduł wraca do sieci biura w ciągu minuty, log wraca. Jeśli moduł zażąda hasła sieci, wpisuje je użytkownik \\
\end{longtable}

\textbf{Ryzyko.} Jeśli w 1d Bluetooth zawiedzie, moduł zostaje bez Wi-Fi. Droga powrotu: sieć modułu z telefonu. Dlatego 1d dopiero po udanych 1b i 1f.


\textbf{Wniosek do analizy GAP, niezależnie od wyniku.} Jeśli sieć modułu nie ma hasła, każdy w pobliżu może sterować ładowarką przez stronę WWW.


\textbf{Narzędzie do napisania.} \texttt{narzedzia\textbackslash{}ble\_\allowbreak{}rpc.\allowbreak{}py}: klient poleceń Shelly przez Bluetooth. Ta sama treść JSON co przez Wi-Fi, opakowana w ramkę Bluetooth Shelly.


\subsection{T2. Funkcje sterownika bez internetu}

\textbf{Cel.} Potwierdzić, że sterownik robi swoje, gdy moduł ma tylko sieć lokalną (poziom C).


\begin{longtable}{>{\raggedright\arraybackslash}p{3.00cm}>{\raggedright\arraybackslash}p{5.62cm}>{\raggedright\arraybackslash}p{6.66cm}}
\toprule
\rowcolor{apteallight}\textbf{Krok} & \textbf{Co robimy} & \textbf{Wynik oczekiwany} \\
\midrule\endhead
\bottomrule\endfoot
2a & Symulator: wolne {\symbolfont→} podłączone {\symbolfont→} ładuje & sterownik przechodzi stany „wolna” {\symbolfont→} „auto podłączone” {\symbolfont→} „ładuje”, zamyka stycznik; moduł pokazuje to w polach \\
2b & Limit prądu z \textbf{menu ładowarki}: 8 A & sterownik sam zgłasza 8 A; pole w module aktualizuje się bez naszego zapisu \\
2c & Limit prądu z modułu przy ładowaniu: 12 A & sterownik potwierdza w 1 s; wyświetlacz pokazuje 12 A \\
2d & Godzina po starcie: włączamy ładowarkę, patrzymy na pierwsze ramki czasu & przez około 20 s moduł odpowiada zerową datą, potem poprawną; sprawdzamy, co w tym czasie pokazuje wyświetlacz \\
2e & Stan sieci wysyłany do sterownika: obserwacja przy poziomie C, a potem w T1d przy wyłączonej stacji Wi-Fi & według opisu protokołu Tuya: 2 = brak routera, 3 = router bez chmury, 4 = połączony z chmurą. Moduł wysyła „4” także bez chmury. Zapisujemy, czy kiedykolwiek wysyła 2 albo 3. Od tego zależy, czy ikona na wyświetlaczu coś znaczy \\
2f & Poziom C: zmiana limitu przez MQTT i przez most OCPP & sterownik potwierdza tak jak 25.09; broker i most działają w sieci lokalnej bez internetu \\
\end{longtable}

\subsection{T3. Okno godzin, tryb pracy, limit energii}

\textbf{Cel.} Trzy funkcje, które Tuya obsługuje z aplikacji, a my na razie tylko z menu. Ustalić, czy sterownik je zgłasza i wykonuje, oraz co trzeba dodać w produkcie Shelly.


\textbf{Narzędzie do napisania: serwer czasu z przesunięciem.} \texttt{narzedzia\textbackslash{}zegar\_\allowbreak{}przesuniety.\allowbreak{}py} podaje czas przesunięty tak, żeby za 2 minuty wybiła pełna godzina. Okno godzin ma tylko pełne godziny, więc bez tego test trwałby do dwóch godzin. Moduł przekazuje sterownikowi tylko to, co dostał od serwera czasu, więc przesuwamy zegar ładowarki bez ruszania zegara laptopa. Na czas testu zatrzymujemy kontener z serwerem czasu (port 123), a po zmianie czasu restartujemy sam moduł, żeby od razu pobrał nowy.


\begin{longtable}{>{\raggedright\arraybackslash}p{3.00cm}>{\raggedright\arraybackslash}p{6.14cm}>{\raggedright\arraybackslash}p{6.14cm}}
\toprule
\rowcolor{apteallight}\textbf{Krok} & \textbf{Co robimy} & \textbf{Wynik oczekiwany} \\
\midrule\endhead
\bottomrule\endfoot
3a & Serwer z przesunięciem: za 2 min wybije godzina G. Z menu: tryb „harmonogram”, okno od G do G+1. Symulator w stanie „ładuje”. Potem przesunięcie na 2 min przed G+1 & sterownik zgłasza tryb „harmonogram” i okno (G, G+1); przed G stan „czeka”, o G ładowanie rusza, o G+1 staje \\
3b & To samo bez ważnego czasu: serwer czasu zatrzymany, ładowarka włączona na nowo & sterownik dostaje zerową datę; zapisujemy, co robi z oknem: nie startuje, startuje od razu albo coś innego. Potem serwer czasu z powrotem \\
3c & Z menu: tryb „energia”, limit 1 kWh & sterownik zgłasza tryb „energia” i, oby, limit energii; bez obciążenia licznik nie ruszy, więc sprawdzamy tylko zgłaszanie \\
3d & Z menu: sprawdzamy, czy da się ustawić okno z minutami & jeśli tak, założenie „tylko pełne godziny” jest błędne i trzeba je poprawić w zestawieniu punktów danych \\
\end{longtable}

\textbf{Zapis tych trzech ustawień z modułu to osobny etap.} Moduł nie ma polecenia do zapisu dowolnego punktu danych. Zapisuje tylko skrypt usługi, i tylko do pól zdefiniowanych w produkcie. Trzeba więc dodać pola w portalu i wgrać nową konfigurację, a to kasuje Wi-Fi modułu.


\textbf{Wynik uboczny.} Lista pól do dodania w portalu: tryb pracy, okno godzin, limit energii, błędy, sygnał z auta, temperatura, wersja sterownika, energia ostatniej sesji.


\subsection{T4. Harmonogram, webhook i pamięć w module}

\textbf{Cel.} To, co moduł Shelly daje ponad Tuya, ma działać bez chmury i przetrwać restart.


\begin{longtable}{>{\raggedright\arraybackslash}p{3.00cm}>{\raggedright\arraybackslash}p{6.14cm}>{\raggedright\arraybackslash}p{6.14cm}}
\toprule
\rowcolor{apteallight}\textbf{Krok} & \textbf{Co robimy} & \textbf{Wynik oczekiwany} \\
\midrule\endhead
\bottomrule\endfoot
4a & Harmonogram w module: o +2 min ustaw limit 6 A, o +4 min 16 A & oba zadania odpalają o czasie, sterownik potwierdza \\
4b & Najpierw lista zdarzeń, na które moduł umie wywołać webhook (polecenie, które nic nie zmienia) & wiemy, czy zmiana naszych pól (stan, limit, włącznik) jest zdarzeniem, czy tylko zdarzenia systemowe \\
4c & Webhook: przy zmianie stanu ładowania moduł wywołuje adres na laptopie (mały nasłuch HTTP w Pythonie) & wywołanie przy każdej zmianie, z treścią zdarzenia; jeśli 4b pokaże brak takiego zdarzenia, zapisujemy \leavevmode{\symbolfont\color{apred}\textbf{✘}} i powód \\
4d & Pamięć klucz-wartość: zapis wartości testowej & odczyt daje to samo; po T5 sprawdzamy, czy przetrwała zanik zasilania \\
4e & Harmonogram i webhook po restarcie modułu (samo polecenie restartu, bez zaniku zasilania) & nadal są i działają \\
\end{longtable}

\subsection{T5. Zanik zasilania i restart modułu}

\textbf{Cel.} Rozstrzygnąć, co pamięta sterownik, co moduł i kto komu narzuca wartość po starcie.


\textbf{Co przewiduje kod naszego skryptu.} Po starcie skrypt czyta limit i włącznik ze sterownika i nadpisuje nimi pola w module. Do sterownika pisze tylko wtedy, gdy pole się zmieni na wartość inną niż ostatnio otrzymana ze sterownika. Normalnie wygrywa więc sterownik.


Słabe miejsce: zanim sterownik pierwszy raz zgłosi limit, skrypt nie zna jego wartości. Każda zmiana pola w tych kilku sekundach pójdzie do sterownika. Taką zmianą może być odtworzenie zapamiętanej wartości, harmonogram albo polecenie z aplikacji. Dlatego w 5d patrzymy na pierwsze sekundy logu.


\begin{longtable}{>{\raggedright\arraybackslash}p{3.00cm}>{\raggedright\arraybackslash}p{6.14cm}>{\raggedright\arraybackslash}p{6.14cm}}
\toprule
\rowcolor{apteallight}\textbf{Krok} & \textbf{Co robimy} & \textbf{Wynik oczekiwany / co zapisujemy} \\
\midrule\endhead
\bottomrule\endfoot
5a & Stan wyjściowy, zapisany z logu: limit 10 A ustawiony \textbf{z modułu}, ładowanie wyłączone, symulator „podłączone”; w drugiej rundzie dodatkowo tryb „harmonogram” z oknem, limit energii 1 kWh, wpis klucz-wartość i harmonogram z T4 & pełny zrzut wszystkich punktów danych i pól modułu \\
5b & Wyłącznik ładowarki: wyłączyć na 30 s, włączyć & moduł i sterownik startują; monitor wraca \\
5c & Zrzut po starcie (odpowiedź sterownika na zapytanie o wszystkie punkty) porównany z 5a & które wartości sterownik zachował; oczekujemy limitu prądu, trybu, okna i limitu energii; zapisujemy każdą różnicę \\
5d & Pierwsze sekundy logu po starcie & \textbf{czy moduł wysłał zapis limitu albo włącznika, zanim sterownik zgłosił swoje wartości.} Jeśli tak, to jest scenariusz z przykładu poniżej \\
5e & Włącznik ładowania: przed wyłączeniem „włączone”, symulator „ładuje” & czy po powrocie zasilania ładowanie rusza samo, czy czeka na polecenie; to zachowanie sterownika, ważne dla klienta \\
5f & Czas: po ilu sekundach od startu sterownik dostał poprawną godzinę & z logu; wcześniej zerowa data \\
5g & Moduł: konfiguracja Wi-Fi, MQTT, serwer czasu, harmonogram, webhook, wpis klucz-wartość & wszystko na miejscu \\
5h & Energia i czas sesji w polach modułu & energia sesji \textbf{zawsze 0}, bo liczymy ją z licznika całkowitego, którego sterownik nie zgłasza; czas sesji 0, bo trzymany tylko w pamięci RAM \\
5i & Powtórka 5a–5d z limitem ustawionym \textbf{z menu ładowarki} & drugi kierunek: czy moduł po starcie nie nadpisze ustawienia z menu swoją zapamiętaną wartością \\
5j & Restart \textbf{samego modułu} w trakcie ładowania (symulator „ładuje”), sterownik pracuje dalej & czy sterownik przerywa ładowanie, gdy moduł przestaje dawać znak życia; co moduł wysyła po powrocie. Tak będzie przy każdej aktualizacji oprogramowania modułu \\
\end{longtable}

\textbf{Dlaczego to ważne.} Klient ustawia z menu 8 A, bo ma słabe zabezpieczenie. Moduł pamięta z aplikacji 16 A. Po zaniku zasilania moduł wstaje, wysyła 16 A i ładowarka rusza z 16 A. Taki błąd byłby groźny, więc 5d i 5i to najważniejsze punkty planu.


\textbf{Rekomendacja, jeśli 5c pokaże, że sterownik sam pamięta limit i włącznik.} Wyłączyć w portalu zapamiętywanie tych dwóch pól. Wtedy nic nie dają, bo sterownik i tak zgłasza swoje wartości po starcie, a są jedynym źródłem opisanego ryzyka.


\subsection{T6. Reset sieci z menu ładowarki}

\textbf{Cel.} W Q11 na Tuya pozycja menu „reset sieci” wprowadza moduł w parowanie. Sprawdzić, co to polecenie sterownika robi z modułem Shelly.


\begin{longtable}{>{\raggedright\arraybackslash}p{3.00cm}>{\raggedright\arraybackslash}p{5.55cm}>{\raggedright\arraybackslash}p{6.73cm}}
\toprule
\rowcolor{apteallight}\textbf{Krok} & \textbf{Co robimy} & \textbf{Wynik oczekiwany} \\
\midrule\endhead
\bottomrule\endfoot
6a & Monitor na najwyższym poziomie logu; z menu ładowarki: reset sieci & w logu ramka od sterownika z poleceniem resetu Wi-Fi i reakcja modułu: ignoruje, kasuje Wi-Fi albo wchodzi w parowanie \\
6b & Jeśli moduł skasował Wi-Fi: ponowne dodanie do sieci & moduł wraca do sieci \\
\end{longtable}

\textbf{Drogi powrotu w 6b, w tej kolejności:} (1) parowanie z telefonu przez Bluetooth w aplikacji Shelly; 24.09 zawiodło, bo aplikacja zapisała nazwę sieci ze spacją na końcu. (2) Sieć modułu i jego strona WWW: nazwa sieci i hasło wpisane przez użytkownika. USB odpada, bo moduł jest zasilany z płytki.


\textbf{Ryzyko.} Utrata Wi-Fi modułu, dlatego ten test jest na końcu. Wynik zapisujemy jako fakt do dokumentacji dla fabryki.


\subsection{T7. API}

\textbf{Cel.} Zobaczyć, jak działa API modułu i gdzie są jego granice. W module jest 166 poleceń w 30 grupach. Sprawdzamy te, które mają znaczenie dla klienta i dla naszego zaplecza.


\begin{longtable}{>{\raggedright\arraybackslash}p{3.00cm}>{\raggedright\arraybackslash}p{5.86cm}>{\raggedright\arraybackslash}p{6.42cm}}
\toprule
\rowcolor{apteallight}\textbf{Krok} & \textbf{Co robimy} & \textbf{Wynik oczekiwany / co zapisujemy} \\
\midrule\endhead
\bottomrule\endfoot
7a & Wartości spoza zakresu przez API: limit 5, 17, 20 i 32 A; tekst zamiast liczby & kto odrzuca: pole modułu, skrypt czy sterownik; co trafia do sterownika. Definicja Tuya dopuszcza 6–32 A, a Q11 ma 16 A. Z symulatorem bez obciążenia to bezpieczne \\
7b & Wychodzący WebSocket: moduł sam łączy się z serwerem testowym na laptopie & co moduł wysyła sam; czy serwer może odesłać polecenie, np. limit 10 A, i czy dotrze ono do sterownika. To podstawa tłumacza OCPP na serwerze AmperePoint \\
7c & Połączenia naraz: most, monitor, strona WWW i kolejne klienty WebSocket & przy którym połączeniu moduł zaczyna odmawiać albo zwalniać \\
7d & Kopia konfiguracji na laptop & co zawiera; jeśli hasło Wi-Fi, plik nie może trafić do repozytorium \\
7e & Hasło do API: ustawiamy, sprawdzamy HTTP, WebSocket, Bluetooth i stronę WWW; potem zdejmujemy & które drogi wymagają hasła. Most OCPP i monitor przestaną działać, dopóki nie dopiszemy logowania. Hasło zapisujemy w pliku projektu poza repozytorium, bo zgubione oznacza reset fabryczny \\
7f & Opcjonalnie: kanał UDP & czy działa i do czego mógłby służyć \\
7g & Home Assistant: dodanie modułu przez oficjalną integrację Shelly & czy HA widzi pola ładowarki i czy może nimi sterować; lokalnie, bez chmury. Wymaga dodania urządzenia w naszym HA \\
\end{longtable}

\section{Kolejność i czas}

\begin{longtable}{>{\raggedright\arraybackslash}p{3.00cm}>{\raggedright\arraybackslash}p{4.04cm}>{\raggedright\arraybackslash}p{1.96cm}>{\raggedright\arraybackslash}p{1.33cm}>{\raggedright\arraybackslash}p{1.82cm}>{\raggedright\arraybackslash}p{1.81cm}}
\toprule
\rowcolor{apteallight}\textbf{Kolejność} & \textbf{Test} & \textbf{Poziom} & \textbf{Ryzyko} & \textbf{Czas} & \textbf{Potrzebny użytkownik} \\
\midrule\endhead
\bottomrule\endfoot
1 & T1a–c, T1f Bluetooth i strona WWW & D przy działającym Wi-Fi & małe & 1 h 15 min z narzędziem & tak, telefon \\
2 & Przejście modułu na poziom C & C & małe & 15 min & nie \\
3 & T1g aplikacja w sieci lokalnej, T2 funkcje sterownika & C, E & małe & 55 min & tak, telefon, menu i symulator \\
4 & T5 zanik zasilania, pierwsza runda, z 5j & C & małe & 1 h & tak, wyłącznik \\
5 & T3 okno, tryb, limit energii & C, E & małe & 1 h z narzędziem & tak, menu \\
6 & T4 harmonogram, webhook, pamięć & C & małe & 45 min & nie \\
7 & T7a–d i 7g API, Home Assistant & C & małe & 1 h 15 min & nie \\
8 & T5 druga runda, po T3 i T4 & C & małe & 30 min & tak \\
9 & T7e hasło do API & C & średnie & 20 min & nie \\
10 & T1d–e Bluetooth bez Wi-Fi & D & średnie & 30 min & tak, telefon jako droga awaryjna \\
11 & T6 reset sieci z menu & E & duże & 20 min plus ewentualne dodanie do sieci & tak \\
12 & Powrót modułu na adres automatyczny, poziom B & B & małe & 10 min & nie \\
\end{longtable}

Razem około 8 godzin, czyli dwa dni pracy. Po wszystkim: tabela wyników do analizy GAP v4 i uzupełnienie dokumentacji dla fabryki o wyniki T5 i T6.


\section{Czego ten plan nie obejmuje}

\begin{itemize}
\item Prąd i moc faz pod obciążeniem, pełna sesja z energią: wymaga obciążenia, osobny test.
\item Zapis trybu, okna godzin i limitu energii z modułu: wymaga nowych pól w produkcie, osobny etap po T3.
\item Aktualizacja sterownika przez moduł: wymaga pliku od fabryki.
\item Ekran ładowarki w aplikacji Shelly: sprawa portalu, nie łączności.
\item Porównanie z Tuya bez internetu (T0): tylko jeśli administrator sieci odetnie jedną ładowarkę od internetu.
\end{itemize}

\end{document}